\FloatBarrier
\section{Data and Training Details}\label{app:data}

\FloatBarrier
\subsection{Data}
We use the publicly available National Vulnerability Database (NVD)\footnote{\url{https://nvd.nist.gov}} for vulnerability data.
Up until recently the NVD enriched nearly all CVEs with CWEs, CVSS base scores, severity labels and more, but since April 2026~\citep{nist2026nvd} it only enriches vulnerabilities it considers high-priority.

We include all CVEs with a CWE, primary\footnote{When a CVE has several CVSS entries, we use the primary score assigned by NIST.} base CVSS v3 score and severity, published from 2022 to 2024 (both included), for a total of \nAll.
Of these, \nVThreeZero{} have CVSS v3.0 scores, and the rest v3.1.
The formula for computing the base score from the eight components gives the same score in both versions for every one of these CVEs.

For training/validation/test we use a temporal split, to prevent data leakage and better reflect real use: We train and validate on older CVEs, and test on newer ones.
The \nTrainVal{} CVEs with identifiers from 2022 and 2023 are split into training and validation data randomly with an 80/20 ratio, stratified by severity, resulting in \nTrain{} training and \nVal{} validation CVEs, Table~\ref{tab:data}.
The test data are the \nTest{} remaining (i.e., usable) CVEs, all with identifiers from 2024.
Note that Low severity is less than 2\% of all CVEs, in all splits.
The CWEs exhibit a long tail, Table~\ref{tab:data}.
There are \nTrainCwes{} different CWEs in the training set and \nTestCwes{} in the test set, but the top 100 CWEs in the training set occur in \nTopCov\% of the test CVEs.
We do not use the recent VLAI data and model of \citet{bonhomme2025vlai} because the random train/validation/test split of the underlying dataset, \texttt{CIRCL/vulnerability-scores}, includes CVEs from 2024 (also present in our test data), 2025 and 2026, so VLAI's performance could in part reflect memorisation rather than generalisation ability.
A temporal split is necessary to fairly compare all systems' abilities to generalise, particularly since products mentioned in descriptions repeat often.

\FloatBarrier
\subsection{TF-IDF Baselines}
As baselines we implement TF-IDF classification/regression methods similar to previous work~\citep{elbaz2020fighting}.
We create TF-IDF representations with unigrams and bigrams, using sublinear tf and min document frequency 2, including stop words since words such as ``not'' or ``without'' can affect vulnerability meaning.
We limit the number of features to 50,000.
We train a different model for each output, Table~\ref{tab:summary}, using scikit-learn~\citep{pedregosa2011sklearn}:
\begin{description}\sloppy
  \item[LogReg] We use logistic regression, and select the regularisation strength, $C \in \{0.3, 1, 3, 10\}$, based on validation set accuracy, separately for severity and CWE classification, and for base component classification when predicting all eight components at once.
  \item[CB-LogReg] Class-balanced logistic regression: Identical setup, except we weight each example inversely proportional to its true class frequency.
  \item[Rand.\ Forest] Untuned random forest classifier, with 300 trees for severity and CWE, or 100 for each base component.
  \item[Ridge] Untuned ridge regression to predict scores, except that we choose the $\alpha \in \{0.3, 1, 3, 10\}$ with lowest MAE in validation.
\end{description}
For CWE we only train the models on the most frequent 100 CWEs in the training data, plus an \textsc{Other} class for the remaining CWEs.
A prediction of \textsc{Other} is always considered an error since we compare with the CVEs' true CWEs.

\FloatBarrier
\subsection{Language Model}
We use the 3.8B-parameter phi-3.5-mini-instruct (Phi) causal LM (decoder-only LLM)~\citep{abdin2024phi3}.
Unlike VulnWatch~\citep{databricks2025vulnwatch}, a system that utilises hosted (commercial) LLMs, Phi is available for download, and specialised enough, to be practically usable on local hardware.
Phi's training data is reported to end in October 2023,\footnote{\url{https://huggingface.co/microsoft/Phi-3.5-mini-instruct}} i.e., before any CVEs in the test set.

We fine-tune with LoRA~\citep{hu2022lora} using MLX~\citep{hannun2023mlx}, on a single Apple M5 Pro chip with 64\,GB unified memory.
The LoRAs (rank 16, scale 20, dropout 0.05) are applied to all linear projections in the final 16 transformer layers, totalling \nLoraParams M parameters, only \nLoraPct\% of Phi's parameters.
We train for 5,000 steps (one epoch) with batch size 8.
The learning rate is $10^{-4}$ for the first 100 (warm-up) steps, then decays cosinally to $10^{-5}$.
We exclude \nDropped{} examples (\nDroppedPct\%) with more than 768 tokens to limit memory use.
In contrast we include all test set CVEs, irrespective of length.

The fine-tuning target is a single JSON object: First the eight base score components, then CWE, score, and severity.
Placing score after components ensures that Phi's output is internally consistent, since it has already (essentially)\footnote{Generation is not constrained, so the model can still write a score that disagrees with its vector. It does so on less than 1\% of test CVEs (Table~\ref{tab:consistency}).} committed to a vector by the time it outputs the score, similar to our formula scoring method (Section~\ref{sec:evaluation}).
We evaluate a saved model every 1,000 training steps on a sample of 1,000 validation CVEs, scoring predicted scores via formula, and keep the one with best severity accuracy using greedy decoding, Table~\ref{tab:ckpt}.

For zero-shot Phi we use a prompt with a brief task description, all output field names and the legal values for components and severity (Appendix~\ref{app:prompts}).
Zero-shot Phi is very verbose and slow, so to limit run time we only include it on a random sample of \nSub{} CVEs from the test set.
Where relevant, we evaluate other systems on the same sample.

\FloatBarrier
\subsection{Dataset Statistics}\label{app:stats}
\begin{table}[t]
\centering
\caption{\textbf{Data splits.} Severity columns give the share of CVEs (\%) in each tier.}
\label{tab:data}
\small
\begin{tabular}{lrrrrrr}
\toprule
Split & CVEs & Critical & High & Medium & Low & CWEs \\
\midrule
Train (2022--23) & 40{,}436 & 15.4 & 39.3 & 43.7 & 1.7 & 464 \\
Validation (2022--23) & 10{,}109 & 15.4 & 39.3 & 43.7 & 1.7 & 293 \\
Test (2024) & 36{,}325 & 13.2 & 36.2 & 48.6 & 1.9 & 512 \\
\bottomrule
\end{tabular}
\end{table}

